\documentclass{article}

% Ćirilica: kodiranje fonta i jezik (srpski, ćirilica)
\usepackage[T2A]{fontenc}
\usepackage[utf8]{inputenc}
\usepackage[serbianc]{babel}

% Veličina strane i margine
\usepackage[a4paper,top=2cm,bottom=2cm,left=3cm,right=3cm,marginparwidth=1.75cm]{geometry}

% Korisni paketi
\usepackage{amsmath}
\usepackage{graphicx}
\usepackage{float}
\usepackage{booktabs}
\usepackage{tcolorbox}
\tcbuselibrary{breakable}
\usepackage[colorlinks=true, allcolors=blue]{hyperref}

% Цртежи (TikZ)
\usepackage{tikz}
\usetikzlibrary{arrows.meta, calc, decorations.markings, babel}

% TikZ стилови који се понављају на цртежима
\tikzset{
    naboj/.style={circle, draw, thick, minimum size=6mm, inner sep=0pt, font=\small},
    jezgro/.style={circle, draw, thick, fill=gray!15, minimum size=7mm, inner sep=0pt, font=\small},
    elektron/.style={circle, draw, fill=cyan!30, inner sep=0pt, minimum size=5pt},
    supljina/.style={circle, draw, thick, minimum size=5mm, inner sep=0pt, font=\small},
    sredstrelica/.style={postaction={decorate}, decoration={markings, mark=at position 0.55 with {\arrow{Stealth}}}},
    % --- графици наизменичних величина ---
    pomocna/.style={dashed, gray!70, thin},
    mreza/.style={gray!25, very thin},
    kriva/.style={very thick, red!75!black},
    % --- проводник у пресеку: празан, тачка (ка нама), крстић (од нас) ---
    provodnik/.style={circle, draw, thick, fill=white, minimum size=8pt, inner sep=0pt},
    provodnikT/.style={provodnik, path picture={\fill[black] (path picture bounding box.center) circle (1.1pt);}},
    provodnikK/.style={provodnik, path picture={
        \draw (path picture bounding box.south west) -- (path picture bounding box.north east)
              (path picture bounding box.north west) -- (path picture bounding box.south east);}},
    oznaka/.style={fill=white, inner sep=1pt, font=\small},
}

% Оквир за решене примере
\newtcolorbox[auto counter]{primer}{breakable,
    colback=white, colframe=black!55, colbacktitle=black!55, coltitle=white,
    fonttitle=\bfseries\small, title={Пример~\thetcbcounter},
    boxrule=0.5pt, arc=1.5pt, left=6pt, right=6pt, top=3pt, bottom=3pt,
    before skip=8pt, after skip=10pt}

% Оквир за кључне формуле
\newtcolorbox{kljucno}{colback=blue!4, colframe=blue!45!black,
    boxrule=0.5pt, arc=1.5pt, left=8pt, right=8pt, top=4pt, bottom=4pt}

\title{ОЕТ 2 - Уводно предавање}
\author{Саобраћајно техничка школа, Београд, Земун}

\begin{document}
\maketitle

%==================================================================
% БЕЛЕШКЕ ИЗ СВЕСКЕ (редом по сликама IMG_6501 -- IMG_6508)
%==================================================================

\section{Струја у колу}

\subsection{Једносмерна и наизменична струја}

Струја је усмерено кретање наелектрисања. Према томе како се мења током времена, делимо је на \textbf{једносмерну} и \textbf{наизменичну}.

\textbf{Једносмерна струја (DC)} тече увек у истом смеру. Најчешће јој је и јачина стална. Добијамо је из батерије, акумулатора или исправљача (пуњача).

\textbf{Наизменична струја (AC)} периодично мења и смер и јачину. Најчешће има облик синусоиде. Таква је струја у градској мрежи.

\begin{figure}[H]
    \centering
    \begin{tikzpicture}[>=Stealth]
        % --- (а) једносмерна ---
        \draw[->] (0,0) -- (4.2,0) node[below] {$t$};
        \draw[->] (0,-1.4) -- (0,1.6) node[left] {$i$};
        \draw[very thick, blue] (0,0.9) -- (3.9,0.9);
        \node[left] at (0,0.9) {$I$};
        \node at (2,-1.9) {(а) једносмерна};
        % --- (б) наизменична ---
        \begin{scope}[xshift=5.6cm]
            \draw[->] (0,0) -- (4.2,0) node[below] {$t$};
            \draw[->] (0,-1.4) -- (0,1.6) node[left] {$i$};
            \draw[very thick, red, domain=0:3.9, samples=120] plot (\x, {sin(180*\x/1.5)});
            % период T: између два суседна максимума (x = 0,75 и x = 3,75)
            \draw[dashed] (0.75,1) -- (0.75,1.3);
            \draw[dashed] (3.75,1) -- (3.75,1.3);
            \draw[<->] (0.75,1.2) -- (3.75,1.2) node[midway, above] {$T$};
            \node at (2,-1.9) {(б) наизменична};
        \end{scope}
    \end{tikzpicture}
    \caption{Струја током времена: (а) стално у истом смеру, (б) мења смер}
    \label{fig:dc-ac}
\end{figure}

Наизменичну струју описујемо са две величине:
\begin{itemize}
    \item \textbf{Период} $T$ -- време за које се заврши један цео циклус. Јединица је секунда (s).
    \item \textbf{Фреквенција} $f$ -- колико се циклуса деси у једној секунди. Јединица је херц (Hz).
\end{itemize}

\begin{equation}
    f = \frac{1}{T}
\end{equation}

У градској мрежи напон је 230~V, а фреквенција $f = 50$~Hz. Значи да је $T = \frac{1}{50}~\mathrm{s} = 20$~ms и да струја промени смер 100 пута у секунди.

Код наизменичне струје електрони не путују дуж жице. Они само осцилују напред-назад, а енергија се ипак преноси до потрошача.

%==================================================================
% ДРУГИ ДЕО: ПАРАМЕТРИ НАИЗМЕНИЧНИХ ВЕЛИЧИНА (уџбеник, стр. 12--21)
%==================================================================

\section{Параметри наизменичних величина}

Величине које се мењају по синусном закону зову се \textbf{простопериодичне}. Тако се мењају струја, напон и електромоторна сила.

Наизменичне величине обележавамо малим словима ($i$, $u$, $e$). Сталне величине обележавамо великим словима ($I_m$, $T$, \dots).

Наизменичну струју потпуно описују сталне величине -- \textbf{параметри}: 
\begin{itemize}
\item тренутна и максимална вредност
\item периода
\item фреквенција
\item кружна фреквенција
\item фаза и почетна фаза
\item средња вредност
\item ефективна вредност
\end{itemize}

\subsection{Тренутна и максимална вредност}

Струја у било ком тренутку $t$ је:
\begin{equation}
    i = I_m\sin\omega t
\end{equation}
\begin{itemize}
    \item $i$ је \textbf{тренутна вредност} -- вредност струје у једном тренутку. У тренутку $t_1$ струја је $i_1 = I_m\sin\omega t_1$, у тренутку $t_2$ је $i_2 = I_m\sin\omega t_2$.
    \item $I_m$ је \textbf{максимална вредност} или \textbf{амплитуда} -- највећа вредност коју струја достиже.
\end{itemize}

\begin{figure}[H]
    \centering
    \begin{tikzpicture}[>=Stealth, font=\small]
        \def\A{1.5}
        \draw[->] (0,0) -- (9,0) node[right] {$t$};
        \draw[->] (0,-1.9) -- (0,2.0) node[left] {$i$};
        \node[below left, font=\footnotesize] at (0,0) {$0$};
        \draw[pomocna] (0,\A) -- (8.4,\A);
        \draw[pomocna] (0,-\A) -- (8.4,-\A);
        \node[left] at (0,-\A) {$-I_m$};
        \foreach \x/\l in {2/{\dfrac{T}{4}},4/{\dfrac{T}{2}},6/{\dfrac{3T}{4}},8/{T}} {
            \draw[pomocna] (\x,-\A) -- (\x,\A);
            \node[below, font=\footnotesize] at (\x,-\A) {$\l$};
        }
        \draw[kriva, domain=0:8.6, samples=100] plot (\x,{\A*sin(45*\x)});
        \node[red!75!black, above right] at (3.0,{\A*sin(135)}) {$i$};
        % амплитуда
        \draw[<->, thick] (2,0) -- (2,\A) node[midway, right] {$I_m$};
        % тренутне вредности
        \foreach \t/\n in {0.6/1, 1.2/2} {
            \draw[densely dashed, gray] (\t,0) -- (\t,{\A*sin(45*\t)}) -- (0,{\A*sin(45*\t)});
            \fill (\t,{\A*sin(45*\t)}) circle (1.4pt);
            \node[below, font=\footnotesize] at (\t,0) {$t_\n$};
            \node[left, font=\footnotesize] at (0,{\A*sin(45*\t)}) {$i_\n$};
        }
    \end{tikzpicture}
    \caption{Тренутне вредности $i_1$ и $i_2$ и амплитуда $I_m$}
    \label{fig:trenutna}
\end{figure}

\subsection{Периода}

Наизменична струја је периодична: после одређеног времена све њене вредности се понављају.

\textbf{Периода} $T$ је време за које се изврши једна цела промена. Јединица је секунда (s).
\begin{itemize}
    \item Периоду је најлакше очитати између два суседна максимума или два суседна минимума.
    \item Половина периоде зове се \textbf{полупериода} ($T/2$).
    \item Током једне периоде струја достиже један максимум и један минимум.
\end{itemize}

\begin{figure}[H]
    \centering
    \begin{tikzpicture}[>=Stealth, font=\small]
        \def\A{1.3}
        \def\s{0.5}  % cm по ms
        \foreach \t in {1,...,18}
            \draw[mreza] (\t*\s,-\A) -- (\t*\s,\A);
        \draw[pomocna] (0,\A) -- (18.3*\s,\A);
        \draw[pomocna] (0,-\A) -- (18.3*\s,-\A);
        \draw[->] (0,0) -- (19.3*\s,0) node[right] {$t\,[\mathrm{ms}]$};
        \draw[->] (0,-1.8) -- (0,1.95) node[left] {$i$};
        \foreach \t in {2,4,...,18}
            \node[below, font=\scriptsize] at (\t*\s,0) {\t};
        \draw[kriva, domain=0:9.25, samples=150] plot (\x,{\A*cos(72*\x-72)});
        % T између два максимума
        \draw[pomocna] (1,\A) -- (1,1.85);
        \draw[pomocna] (6,\A) -- (6,1.85);
        \draw[<->] (1,1.7) -- (6,1.7) node[midway, above] {$T$};
        % T између два минимума
        \draw[pomocna] (3.5,-\A) -- (3.5,-1.85);
        \draw[pomocna] (8.5,-\A) -- (8.5,-1.85);
        \draw[<->] (3.5,-1.7) -- (8.5,-1.7) node[midway, below] {$T$};
    \end{tikzpicture}
    \caption{Очитавање периоде: максимуми су у 2~ms и 12~ms, минимуми у 7~ms и 17~ms, па је $T = 10$~ms}
    \label{fig:perioda}
\end{figure}

\subsection{Фреквенција}

\textbf{Фреквенција} (учестаност) $f$ је број периода у једној секунди:
\begin{equation}
    f = \frac{1}{T}
\end{equation}
Јединица је херц (Hz), по немачком физичару Хајнриху Херцу (1857--1894):
\[
    1~\mathrm{Hz} = \frac{1}{1~\mathrm{s}}
\]
Периода и фреквенција су обрнуто сразмерне. Што је периода краћа, фреквенција је већа.

\subsection{Кружна фреквенција}

\textbf{Кружна фреквенција} (кружна учестаност) $\omega$ је угаона брзина навојка у генератору.

За једну периоду $T$ навојак пређе пун круг: $360^\circ$, односно $2\pi$ радијана. Зато је $\omega\cdot T = 2\pi$, па је:
\begin{equation}
    \omega = \frac{2\pi}{T} = 2\pi f
\end{equation}
Јединица је радијан у секунди (rad/s).

Табела повезује време, угао у степенима и угао у радијанима:
\begin{table}[H]
    \centering
    \begin{tabular}{l|ccccc}
        \toprule
        време $t$                        & $0$ & $\frac{T}{4}$   & $\frac{T}{2}$ & $\frac{3T}{4}$   & $T$ \\
        угао $\omega t$ (степени)        & $0^\circ$ & $90^\circ$ & $180^\circ$ & $270^\circ$ & $360^\circ$ \\
        угао $\omega t$ (радијани)       & $0$ & $\frac{\pi}{2}$ & $\pi$         & $\frac{3\pi}{2}$ & $2\pi$ \\
        \bottomrule
    \end{tabular}
    \caption{Веза времена и угла током једне периоде}
    \label{tab:ugao}
\end{table}

\begin{primer}
Фреквенција мрежног напона је $f = 50$~Hz. Колика је периода, а колика кружна фреквенција?

\textbf{Решење.}
\[
    T = \frac{1}{f} = \frac{1}{50}~\mathrm{s} = 0{,}02~\mathrm{s} = 20~\mathrm{ms}
\]
\[
    \omega = 2\pi f = 2\cdot 3{,}14\cdot 50 = 314~\frac{\mathrm{rad}}{\mathrm{s}}
    \qquad\text{или}\qquad
    \omega = \frac{2\pi}{T} = \frac{2\pi}{20\cdot 10^{-3}~\mathrm{s}} = 314~\frac{\mathrm{rad}}{\mathrm{s}}
\]
\end{primer}

\subsection{Фаза и почетна фаза}

Посматрање струје може почети у било ком тренутку. Тај тренутак обележавамо са $t = 0$. Струја тада у општем случају није нула:
\begin{equation}
    i = I_m\sin(\omega t + \psi)
\end{equation}
\begin{itemize}
    \item $\omega t + \psi$ је \textbf{фаза} (фазни угао). Она показује у ком делу периоде се струја налази.
    \item $\psi$ је \textbf{почетна фаза} -- фаза у тренутку $t = 0$. Тада је $i(0) = I_m\sin\psi$.
\end{itemize}
Ако је $\psi > 0$, синусоида је померена \textbf{улево}: кроз нулу пролази пре $t = 0$.
Ако је $\psi < 0$, синусоида је померена \textbf{удесно}: кроз нулу пролази после $t = 0$.

\begin{figure}[H]
    \centering
    \begin{tikzpicture}[>=Stealth, font=\small]
        \def\A{1.4}
        \def\k{40}  % степени по cm
        \draw[->] (-2.6,0) -- (9.6,0) node[right] {$\omega t$};
        \draw[->] (0,-1.8) -- (0,2.1) node[above] {$i$};
        \draw[pomocna] (0,\A) -- (9.2,\A);
        \draw[pomocna] (0,-\A) -- (9.2,-\A);
        \foreach \d/\r in {90/{\dfrac{\pi}{2}},180/{\pi},270/{\dfrac{3\pi}{2}},360/{2\pi}} {
            \draw[pomocna] ({\d/\k},-\A) -- ({\d/\k},\A);
            \node[below, font=\footnotesize, align=center] at ({\d/\k},-\A) {$\d^\circ$\\[1pt]$\r$};
        }
        % ψ = 0 (референтна) и ψ = 60°
        \draw[thick, gray, dashed, domain=-100:370, samples=120] plot ({\x/\k},{\A*sin(\x)});
        \draw[kriva, domain=-100:370, samples=120] plot ({\x/\k},{\A*sin(\x+60)});
        % помак psi између максимума
        \draw[pomocna] (0.75,\A) -- (0.75,1.85);
        \draw[pomocna] (2.25,\A) -- (2.25,1.85);
        \draw[<->] (0.75,1.7) -- (2.25,1.7) node[midway, above] {$\psi$};
        % вредност у t = 0
        \fill[red!75!black] (0,{\A*sin(60)}) circle (1.7pt);
        \node[pin={[pin distance=7mm, font=\footnotesize]160:{$i(0)=I_m\sin\psi$}}, inner sep=0pt] at (0,{\A*sin(60)}) {};
        % легенда
        \draw[kriva] (5.6,2.05) -- (6.2,2.05) node[right, black, font=\footnotesize] {$\psi = 60^\circ$};
        \draw[thick, gray, dashed] (7.7,2.05) -- (8.3,2.05) node[right, black, font=\footnotesize] {$\psi = 0$};
    \end{tikzpicture}
    \caption{Струја са почетном фазом $\psi = 60^\circ$ (пуна линија) престиже струју са $\psi = 0$ (испрекидана) за $60^\circ$}
    \label{fig:faza}
\end{figure}

Почетна фаза се са графика одређује из тренутка $t_0$ у коме величина пролази кроз нулу \emph{растући}:
\begin{equation}
    \psi = -\frac{t_0}{T}\cdot 360^\circ
\end{equation}

\begin{primer}
Струја је дата изразом $i = 3\sin\!\left(3{,}14\,t + \frac{\pi}{4}\right)$~A. Одредити њене параметре и вредност у тренутку $t = 2$~s.

\textbf{Решење.} Параметри се читају из израза:
\[
    I_m = 3~\mathrm{A}, \qquad \omega = 3{,}14~\frac{\mathrm{rad}}{\mathrm{s}}, \qquad \psi = \frac{\pi}{4} = 45^\circ
\]
\[
    f = \frac{\omega}{2\pi} = \frac{3{,}14}{6{,}28} = 0{,}5~\mathrm{Hz}, \qquad T = \frac{1}{f} = 2~\mathrm{s}
\]
Вредност у $t = 2$~s:
\[
    i = 3\sin(3{,}14\cdot 2 + 0{,}785) = 3\sin(6{,}28 + 0{,}785)~\mathrm{A} \approx 2{,}1~\mathrm{A}
\]
Пошто је $t = 2$~s тачно једна периода, струја има исту вредност као у $t = 0$: $i(0) = 3\sin 45^\circ \approx 2{,}1$~A.
\end{primer}

\begin{primer}
Написати аналитички облик напона ако је познат његов временски дијаграм.
\begin{center}
    \begin{tikzpicture}[>=Stealth, font=\small]
        \def\s{0.5}   % cm по ms
        \def\v{0.7}   % cm по V
        \foreach \t in {1,...,18}
            \draw[mreza] (\t*\s,-2*\v) -- (\t*\s,2*\v);
        \foreach \u in {-2,-1.5,...,2}
            \draw[mreza] (0,\u*\v) -- (18.3*\s,\u*\v);
        \draw[->] (0,0) -- (19.4*\s,0) node[right] {$t\,[\mathrm{ms}]$};
        \draw[->] (0,-2.2*\v) -- (0,2.6*\v) node[left] {$u\,[\mathrm{V}]$};
        \foreach \u/\l in {0.5/{0,5}, 1/{1,0}, 1.5/{1,5}, 2/{2,0}}
            \node[left, font=\scriptsize] at (0,\u*\v) {\l};
        \foreach \t in {2,4,...,18}
            \node[below, font=\scriptsize] at (\t*\s,-2*\v) {\t};
        \draw[kriva, domain=0:9.2, samples=150] plot (\x,{2*\v*sin(72*\x-144)});
    \end{tikzpicture}
\end{center}

\textbf{Решење.} Са графика очитавамо:
\begin{itemize}
    \item амплитуда: $U_m = 2$~V;
    \item напон пролази кроз нулу растући у 4~ms и у 14~ms, па је $T = 14 - 4 = 10$~ms;
    \item кружна фреквенција: $\omega = \dfrac{2\pi}{T} = \dfrac{2\pi}{10\cdot 10^{-3}~\mathrm{s}} = 628~\dfrac{\mathrm{rad}}{\mathrm{s}}$;
    \item почетна фаза ($t_0 = 4$~ms): $\psi = -\dfrac{4}{10}\cdot 360^\circ = -144^\circ$.
\end{itemize}
\[
    u(t) = 2\sin(628\,t - 144^\circ)~\mathrm{V}
\]
Провера: $u(0) = 2\sin(-144^\circ) \approx -1{,}18$~V, што се и види на графику.
\end{primer}

%==================================================================
% ТРЕЋИ ДЕО: СРЕДЊА И ЕФЕКТИВНА ВРЕДНОСТ (уџбеник, стр. 21--23)
%==================================================================

\section{Анализа наизменичног сигнала}

Синусни сигнал стално мења вредност. У многим применама згодине а понекад је и неопходно израчунати колико он заправо утиче на наше коло, на пример колико снаге сигнал преноси. За те потребе уводимо две вредности: \textbf{средњу} и \textbf{ефективну}. Обе су једносмерне струје које у нечему замењују наизменичну.

\subsection{Средња вредност}

\textbf{Средња вредност} $I_{sr}$ је једносмерна струја која за исто време пренесе исту количину наелектрисања као наизменична.

Током целе периоде средња вредност је нула, јер је позитивна полупериода једнака негативној. Зато је рачунамо за \textbf{једну полупериоду}:
\begin{equation}
    I_{sr} = \frac{2}{\pi}\,I_m \approx 0{,}637\,I_m
\end{equation}

\begin{figure}[H]
    \centering
    \begin{tikzpicture}[>=Stealth, font=\small]
        \def\A{2}
        \def\Isr{1.2732}
        % површине
        \fill[violet!35] plot[domain=1.318:4.682, samples=60] (\x,{\A*sin(30*\x)}) -- cycle;
        \fill[orange!45] (0,0) -- (0,\Isr) -- (1.318,\Isr)
            -- plot[domain=1.318:0, samples=30] (\x,{\A*sin(30*\x)}) -- cycle;
        \fill[orange!45] (6,0) -- (6,\Isr) -- (4.682,\Isr)
            -- plot[domain=4.682:6, samples=30] (\x,{\A*sin(30*\x)}) -- cycle;
        % осе и помоћне линије
        \draw[->] (0,0) -- (6.9,0) node[right] {$\omega t$};
        \draw[->] (0,-0.3) -- (0,2.5) node[left] {$i$};
        \node[below left, font=\footnotesize] at (0,0) {$0$};
        \draw[pomocna] (0,\A) -- (6.2,\A);
        \node[left] at (0,\A) {$I_m$};
        \draw[pomocna] (3,0) -- (3,\A);
        \node[below] at (3,0) {$\dfrac{\pi}{2}$};
        \node[below] at (6,0) {$\pi$};
        \draw[thick, dashed, violet!70!black] (0,\Isr) -- (6.3,\Isr);
        \node[left, violet!70!black] at (0,\Isr) {$I_{sr}$};
        \draw[kriva, domain=0:6, samples=80] plot (\x,{\A*sin(30*\x)});
    \end{tikzpicture}
    \caption{Средња вредност за једну полупериоду: љубичаста површина изнад $I_{sr}$ једнака је наранџастим површинама испод}
    \label{fig:srednja}
\end{figure}

Исти изрази важе за напон и ЕМС: $U_{sr} = \frac{2}{\pi}U_m$, $\;E_{sr} = \frac{2}{\pi}E_m$.

\subsection{Ефективна вредност}

\textbf{Ефективна вредност} $I_{ef}$ је једносмерна струја која на истом отпорнику, за време једне периоде, развије исту количину топлоте као наизменична.

Загревање не зависи од смера струје. Зато греју обе полупериоде, и позитивна и негативна.
\begin{equation}
    I_{ef} = \frac{I_m}{\sqrt{2}} \approx 0{,}707\,I_m
\end{equation}

\begin{figure}[H]
    \centering
    \begin{tikzpicture}[>=Stealth, font=\small]
        \def\A{1.6}
        \draw[->] (0,0) -- (8.7,0) node[right] {$\omega t$};
        \draw[->] (0,-2.0) -- (0,2.1) node[left] {$i$};
        \node[below left, font=\footnotesize] at (0,0) {$0$};
        \draw[pomocna] (0,\A) -- (8.3,\A);
        \node[left] at (0,\A) {$I_m$};
        \draw[pomocna] (0,-\A) -- (8.3,-\A);
        \node[left] at (0,-\A) {$-I_m$};
        \foreach \x/\l in {2/{\dfrac{\pi}{2}},4/{\pi},6/{\dfrac{3\pi}{2}},8/{2\pi}} {
            \draw[pomocna] (\x,-\A) -- (\x,\A);
            \node[below, font=\footnotesize] at (\x,-\A) {$\l$};
        }
        \draw[thick, dashed, blue!70!black] (0,{\A/sqrt(2)}) -- (8.3,{\A/sqrt(2)});
        \node[left, blue!70!black] at (0,{\A/sqrt(2)}) {$I_{ef}$};
        \draw[kriva, domain=0:8, samples=100] plot (\x,{\A*sin(45*\x)});
    \end{tikzpicture}
    \caption{Ефективна вредност наизменичне струје}
    \label{fig:efektivna}
\end{figure}

\begin{itemize}
    \item Ефективну вредност често пишемо великим словом без индекса: $I$, $U$, $E$.
    \item Амперметри и волтметри, аналогни и дигитални, показују \textbf{ефективну} вредност.
\end{itemize}

\paragraph{Зашто баш $\sqrt{2}$?} Математички израз за ефективну вредност: 

\begin{equation}
    I_{ef} = \sqrt{ \frac{1}{T} \int_0^Tx^2 dt } \\
    x = \sin \omega t 
\end{equation}


\subsection{Поређење вредности}

\begin{table}[H]
    \centering
    \renewcommand{\arraystretch}{1.9}
    \begin{tabular}{lccc}
        \toprule
        Вредност & Израз & Део од $U_m$ & Мрежа \\
        \midrule
        максимална   & $U_m$                         & $1$         & 325~V \\
        ефективна    & $U = \dfrac{U_m}{\sqrt{2}}$   & $0{,}707$   & 230~V \\
        средња (за $T/2$) & $U_{sr} = \dfrac{2}{\pi}U_m$ & $0{,}637$ & 207~V \\
        \bottomrule
    \end{tabular}
    \caption{Максимална, ефективна и средња вредност}
    \label{tab:poredjenje}
\end{table}

Увек важи: $U_{sr} < U < U_m$.

%==================================================================
% КРАТАК ПРЕГЛЕД
%==================================================================

\subsection{Како настаје синусоида: обртање навојка}

Кад проводник сече линије магнетне индукције, у њему се индукује електромоторна сила (ЕМС). Највећа је кад их сече под правим углом:
\begin{equation}
    E = B\,l\,v
\end{equation}
где је $B$ магнетна индукција, $l$ дужина проводника, а $v$ брзина проводника.

У навојку се индукује ЕМС кад се кроз њега мења магнетни флукс $\Phi$. Количник $\Delta\Phi/\Delta t$ је брзина промене флукса. За намотај са $N$ навојака:
\begin{equation}
    e = -N\,\frac{\Delta\Phi}{\Delta t}
\end{equation}
Знак минус значи: кад флукс расте, ЕМС је негативна, а кад флукс опада, ЕМС је позитивна.

\paragraph{Генератор.} Навојак се обрће у хомогеном магнетном пољу, између полова N и S, сталном угаоном брзином $\omega$. ЕМС се преко прстенова и четкица одводи до потрошача. Угао за који се навојак окренуо расте равномерно с временом:
\begin{equation}
    \alpha = \omega t
\end{equation}

На слици~\ref{fig:obrtanje} навојак гледамо са стране. Види се као дуж $A$--$X$, а проводници као кружићи. Тачка ($\odot$) значи да ЕМС у проводнику има смер ка нама, а крстић ($\otimes$) од нас.

\begin{figure}[H]
    \centering
    \begin{tikzpicture}[>=Stealth, font=\small]
        \def\R{1.6}      % полупречник кружења проводника
        \def\gx{3.2}     % где почиње график
        \def\sx{0.021}   % cm по степену на графику
        % --- полови магнета ---
        \fill[red!60] (-2.3,2.2) rectangle (2.3,2.8);
        \node[white, font=\bfseries] at (0,2.5) {N};
        \fill[blue!50] (-2.3,-2.8) rectangle (2.3,-2.2);
        \node[white, font=\bfseries] at (0,-2.5) {S};
        % --- линије магнетне индукције (од N ка S) ---
        \foreach \x in {-1.95,-1.3,-0.65,0,0.65,1.3,1.95}
            \draw[gray!35, sredstrelica] (\x,2.2) -- (\x,-2.2);
        \node[right, gray!50!black] at (1.95,-1.7) {$\vec B$};
        % --- путања проводника и положаји навојка ---
        \draw[gray!70] (0,0) circle (\R);
        \foreach \a in {0,30,60,90}
            \draw[dashed, black!60] (\a+180:\R) -- (\a:\R);
        \fill (0,0) circle (1.3pt);
        % --- угао alpha ---
        \draw[->, thin] (0.5,0) arc (0:60:0.5);
        \node at (45:0.74) {$\alpha$};
        % --- смер обртања ---
        \draw[->, thick] (130:1.95) arc (130:170:1.95);
        \node at (150:2.3) {$\omega$};
        % --- ознаке положаја (пре проводника, да их не прекрију) ---
        \node[oznaka, anchor=north west] at ($(0:\R)+(0.08,-0.10)$) {$X$};
        \node[oznaka, anchor=north west] at ($(30:\R)+(0.10,-0.10)$) {$X'$};
        \node[oznaka, anchor=east]       at ($(60:\R)+(-0.17,-0.02)$) {$X''$};
        \node[oznaka, anchor=east]       at ($(90:\R)+(-0.17,0.02)$) {$X'''$};
        \node[oznaka, anchor=north east] at ($(180:\R)+(-0.08,-0.10)$) {$A$};
        \node[oznaka, anchor=north east] at ($(210:\R)+(-0.10,-0.08)$) {$A'$};
        \node[oznaka, anchor=north east] at ($(240:\R)+(-0.10,-0.08)$) {$A''$};
        \node[oznaka, anchor=west]       at ($(270:\R)+(0.17,0)$) {$A'''$};
        % --- проводници ---
        \node[provodnik] at (0:\R) {};
        \node[provodnik] at (180:\R) {};
        \foreach \a in {30,60,90} {
            \node[provodnikT] at (\a:\R) {};
            \node[provodnikK] at (\a+180:\R) {};
        }
        % --- график ---
        \draw[->] (\gx,0) -- ({\gx+7.95},0) node[right] {$\alpha$};
        \draw[->] (\gx,-2.05) -- (\gx,2.1) node[above] {$e,\ \Phi$};
        \node[below left, font=\footnotesize] at (\gx,0) {$0$};
        \draw[pomocna] (\gx,\R) -- ({\gx+360*\sx},\R);
        \draw[pomocna] (\gx,-\R) -- ({\gx+360*\sx},-\R);
        \foreach \d/\r in {90/{\dfrac{\pi}{2}},180/{\pi},270/{\dfrac{3\pi}{2}},360/{2\pi}} {
            \draw[pomocna] ({\gx+\d*\sx},-\R) -- ({\gx+\d*\sx},\R);
            \node[below, font=\footnotesize, align=center] at ({\gx+\d*\sx},-\R) {$\d^\circ$\\[1pt]$\r$};
        }
        % флукс (косинус) и ЕМС (синус)
        \draw[thick, red!75!black, domain=0:360, samples=90]
            plot ({\gx+\x*\sx},{0.7*\R*cos(\x)});
        \draw[very thick, blue!70!black, domain=0:360, samples=90]
            plot ({\gx+\x*\sx},{\R*sin(\x)});
        \node[blue!70!black, right] at ({\gx+150*\sx},{\R*sin(150)}) {$e$};
        \node[red!75!black, above left] at ({\gx+330*\sx},{0.7*\R*cos(330)}) {$\Phi$};
        % --- пројекције: висина проводника X -> тачка на синусоиди ---
        \foreach \a in {30,60,90} {
            \draw[dashed, blue!50] (\a:\R) -- ({\gx+\a*\sx},{\R*sin(\a)});
            \fill[blue!70!black] ({\gx+\a*\sx},{\R*sin(\a)}) circle (1.6pt);
        }
        \fill[blue!70!black] (\gx,0) circle (1.6pt);
        \node[draw, font=\small, inner sep=3pt] at ({\gx+180*\sx},-2.9) {$\alpha = \omega t$};
    \end{tikzpicture}
    \caption{Обртање навојка и настанак синусоиде. Висина проводника $X$ у сваком положају пренета је на график: тачке падају тачно на синусоиду ЕМС $e$. Флукс $\Phi$ се мења по косинусу.}
    \label{fig:obrtanje}
\end{figure}

\paragraph{Зашто баш синусоида?} Проводник се креће по кругу. Линије поља сече само компонента брзине која је нормална на $\vec B$ (слика~\ref{fig:brzina}):
\[
    v_\perp = v\sin\alpha
\]
Зато ЕМС расте и опада по синусу:
\begin{equation}
    e = E_m\sin\alpha = E_m\sin\omega t
\end{equation}
где је $E_m$ највећа (максимална) ЕМС.

\begin{figure}[H]
    \centering
    \begin{tikzpicture}[>=Stealth, font=\small]
        \def\R{1.5}
        \coordinate (P) at (50:\R);
        \draw[->, thick, gray!60!black] (-2.4,1.0) -- (-2.4,-1.0) node[midway, left] {$\vec B$};
        \draw[gray!70] (0,0) circle (\R);
        \draw[dashed, gray] (0,0) -- (\R+0.3,0);
        \draw[dashed, black!60] (0,0) -- (P);
        \fill (0,0) circle (1.3pt);
        \draw[->, thin] (0.55,0) arc (0:50:0.55);
        \node at (25:0.8) {$\alpha$};
        % брзина и њене компоненте
        \draw[dashed, gray] ($(P)+({-1.4*sin(50)},0)$) -- ++(0,{1.4*cos(50)}) -- ($(P)+(0,{1.4*cos(50)})$);
        \draw[->, thick, black!55] (P) -- ++(0,{1.4*cos(50)}) node[right] {$v_\parallel$};
        \draw[->, thick, blue!70!black] (P) -- ++({-1.4*sin(50)},0) node[left] {$v_\perp$};
        \draw[->, very thick, red!75!black] (P) -- ++(140:1.4) node[above left] {$\vec v$};
        \node[provodnikT] at (P) {};
    \end{tikzpicture}
    \caption{Брзина проводника $\vec v$ и њене компоненте. Линије $\vec B$ сече само $v_\perp = v\sin\alpha$.}
    \label{fig:brzina}
\end{figure}

\begin{table}[H]
    \centering
    \begin{tabular}{cccc}
        \toprule
        $\alpha$ & Раван навојка & Флукс $\Phi$ & ЕМС $e$ \\
        \midrule
        $0^\circ$   & нормална на $\vec B$   & $\Phi_m$  & $0$ \\
        $90^\circ$  & паралелна са $\vec B$  & $0$       & $E_m$ \\
        $180^\circ$ & нормална на $\vec B$   & $-\Phi_m$ & $0$ \\
        $270^\circ$ & паралелна са $\vec B$  & $0$       & $-E_m$ \\
        $360^\circ$ & као на почетку         & $\Phi_m$  & $0$ \\
        \bottomrule
    \end{tabular}
    \caption{Флукс и ЕМС у карактеристичним положајима навојка}
    \label{tab:polozaji}
\end{table}

\textbf{Шта треба запамтити:}
\begin{itemize}
    \item Један пун обрт навојка ($360^\circ$) даје једну периоду синусоиде.
    \item ЕМС је највећа кад је флукс нула, јер се флукс тада најбрже мења. Флукс и ЕМС су померени за $90^\circ$ (четвртину периоде).
    \item Што се навојак брже обрће (веће $\omega$), то је периода краћа.
\end{itemize}

%==================================================================
% ЧЕТВРТИ ДЕО: КОМПЛЕКСНИ БРОЈЕВИ У ЕЛЕКТРОТЕХНИЦИ (белешке из свеске)
%==================================================================

\section{Комплексни бројеви у електротехници}

Наизменичне величине су синусоиде. Рачун са синусима је тежак, а са комплексним бројевима много лакши јер је рачуница у комплексном домену апсолвирана. Зато и користимо комплексну анализу у електротехници стално користимо.

\subsection{Синус и косинус из круга}

Тачка $P$ кружи по кругу полупречника 1 (\textbf{јединични круг}) сталном угаоном брзином $\omega$. Угао $\theta$ између полупречника $OP$ и хоризонталне осе расте с временом.
\begin{itemize}
    \item Пројекција тачке на \textbf{вертикалну} осу је $\sin\theta$. На тој оси тачка се креће горе-доле. Кад тај покрет развучемо у времену, добијамо \textbf{синусоиду}.
    \item Пројекција тачке на \textbf{хоризонталну} осу је $\cos\theta$. На тој оси тачка се креће лево-десно. Развучено у времену, то је \textbf{косинусоида}.
\end{itemize}

\begin{figure}[H]
    \centering
    \begin{tikzpicture}[>=Stealth, font=\small]
        \def\R{1.25}       % полупречник круга
        \def\t{50}         % угао тачке P
        \def\gx{2.0}       % почетак синусоиде
        \def\gy{-2.05}     % почетак косинусоиде
        \def\s{0.015}      % cm по степену (360° = 5,4 cm)
        % --- круг ---
        \draw[->] (-1.65,0) -- (1.75,0);
        \draw[->] (0,-1.6) -- (0,1.75);
        \draw[thick] (0,0) circle (\R);
        \draw[very thick, blue!70!black] (0,0) -- (0,{\R*sin(\t)});
        \draw[very thick, red!75!black] (0,0) -- ({\R*cos(\t)},0);
        \draw[thick] (0,0) -- (\t:\R);
        \fill (\t:\R) circle (1.8pt);
        \node[above right] at (\t:\R) {$P$};
        \draw[->] (0.35,0) arc (0:\t:0.35);
        \node at ({\t/2}:0.58) {$\theta$};
        \draw[->, thick] (60:1.5) arc (60:105:1.5);
        \node at (118:1.6) {$\omega$};
        \node[left, blue!70!black, font=\footnotesize] at (0,{\R*sin(\t)/2}) {$\sin\theta$};
        \node[below, red!75!black, font=\footnotesize] at ({\R*cos(\t)/2-0.05},0) {$\cos\theta$};
        % --- синусоида (десно) ---
        \draw[->] (\gx,0) -- ({\gx+5.9},0) node[right] {$\theta$};
        \draw[->] (\gx,-1.6) -- (\gx,1.75);
        \draw[pomocna] (\gx,\R) -- ({\gx+360*\s},\R);
        \draw[pomocna] (\gx,-\R) -- ({\gx+360*\s},-\R);
        \foreach \d/\l in {90/{\dfrac{\pi}{2}},180/{\pi},270/{\dfrac{3\pi}{2}},360/{2\pi}} {
            \draw[pomocna] ({\gx+\d*\s},-\R) -- ({\gx+\d*\s},\R);
            \node[below, font=\footnotesize] at ({\gx+\d*\s},-\R) {$\l$};
        }
        \draw[very thick, blue!70!black, domain=0:360, samples=80] plot ({\gx+\x*\s},{\R*sin(\x)});
        \node[blue!70!black, above] at ({\gx+90*\s},\R) {$\sin\theta$};
        \draw[dashed, blue!50] (\t:\R) -- ({\gx+\t*\s},{\R*sin(\t)});
        \fill[blue!70!black] ({\gx+\t*\s},{\R*sin(\t)}) circle (1.6pt);
        % --- косинусоида (доле, угао расте надоле) ---
        \draw[->] (-1.65,\gy) -- (1.75,\gy);
        \draw[->] (0,\gy) -- (0,{\gy-5.9}) node[below] {$\theta$};
        \draw[pomocna] (\R,\gy) -- (\R,{\gy-360*\s});
        \draw[pomocna] (-\R,\gy) -- (-\R,{\gy-360*\s});
        \foreach \d/\l in {90/{\dfrac{\pi}{2}},180/{\pi},270/{\dfrac{3\pi}{2}},360/{2\pi}} {
            \draw[pomocna] (-\R,{\gy-\d*\s}) -- (\R,{\gy-\d*\s});
            \node[right, font=\footnotesize] at (\R,{\gy-\d*\s}) {$\l$};
        }
        \draw[very thick, red!75!black, domain=0:360, samples=80] plot ({\R*cos(\x)},{\gy-\x*\s});
        \node[left, red!75!black] at (-\R,{\gy-180*\s}) {$\cos\theta$};
        \draw[dashed, red!50] (\t:\R) -- ({\R*cos(\t)},{\gy-\t*\s});
        \fill[red!75!black] ({\R*cos(\t)},{\gy-\t*\s}) circle (1.6pt);
    \end{tikzpicture}
    \caption{Тачка $P$ кружи по јединичном кругу. Пројекција на вертикалну осу даје синусоиду (десно), а пројекција на хоризонталну осу косинусоиду (доле).}
    \label{fig:krug-sin-cos}
\end{figure}

Полупречник је 1 за јединични круг, па по Питагориној теореми важи:
\begin{equation}
    \cos^2\theta + \sin^2\theta = 1
\end{equation}

Угаона брзина је брзина промене угла. Одатле добијамо како угао зависи од времена:
\[
    \omega = \frac{d\theta}{dt}
    \;\Rightarrow\; \\  d\theta = \omega\,dt
    \;\Rightarrow\; \\ \int d\theta = \int \omega\,dt
    \;\Rightarrow\; \\ \theta = \omega t
\]
(узели смо да је $\theta = 0$ у $t = 0$). То је иста веза као код навојка генератора: $\alpha = \omega t$.

\subsection{Имагинарни бројеви}

Неке квадратне једначине немају решење међу реалним бројевима:
\[
    x^2 - 1 = 0 \;\Rightarrow\; x_1 = 1,\; x_2 = -1
\]
\[
    x^2 + 1 = 0 \;\Rightarrow\; x^2 = -1 \;\Rightarrow\; x_{1,2} = \pm\sqrt{-1}
\]
Ниједан реалан број на квадрат не даје негативан број. Зато уводимо \textbf{имагинарну јединицу}:
\begin{equation}
    j = \sqrt{-1}, \qquad j^2 = -1
\end{equation}
У математици се обележава са $i$. У електротехници пишемо $j$, јер $i$ већ означава струју. Друга једначина сада има решења $x_1 = j$ и $x_2 = -j$.

Степени броја $j$ понављају се на свака четири корака:
\[
    j^1 = j, \qquad j^2 = -1, \qquad j^3 = -j, \qquad j^4 = 1
\]

\subsection{Комплексна раван}

Имагинарни бројеви су посебна „димензија“. Реални и имагинарни део не могу се сабрати у један број, па их пишемо одвојено:
\begin{equation}
    \underline{A} = a + jb
\end{equation}
где је $a = \operatorname{Re}\{\underline{A}\}$ \textbf{реални део}, а $b = \operatorname{Im}\{\underline{A}\}$ \textbf{имагинарни део}. Комплексне величине обележавамо подвученим словом. Промењиву $А$ у овој једнакости називамо \textbf{комплексан  број}. Комплексни број цртамо као тачку у равни. Реални део је на хоризонталној оси (Re), а имагинарни на вертикалној (Im). Тако да реалну страну осе можемо мапирати на 
$\cos\theta$ а имагинарну страну на $\sin\theta$.

Сада спајамо раван и круг. Тачка на растојању $A$ од почетка, под углом $\theta$, има пројекције $a = A\cos\theta$ и $b = A\sin\theta$:
\begin{equation}
    \underline{A} = A(\cos\theta + j\sin\theta)
\end{equation}
\begin{itemize}
    \item $A = |\underline{A}| = \sqrt{a^2 + b^2}$ је \textbf{модуо} (дужина);
    \item $\theta$ је \textbf{аргумент} (угао), $\operatorname{tg}\theta = \dfrac{b}{a}$.
\end{itemize}
\textbf{Конјуговано комплексни} број има супротан знак имагинарног дела. У равни је то слика броја $\underline{A}$ у огледалу Re осе:
\[
    \underline{A}^* = a - jb = A(\cos\theta - j\sin\theta)
\]

\begin{figure}[H]
    \centering
    \begin{tikzpicture}[>=Stealth, font=\small]
        \def\A{2.4}
        \def\t{35}
        \coordinate (A) at (\t:\A);
        \coordinate (K) at (-\t:\A);
        \draw[->] (-0.6,0) -- (3.4,0) node[right] {Re};
        \draw[->] (0,-1.9) -- (0,2.0) node[above] {Im};
        \draw[gray!50, dotted, thick] (-55:\A) arc (-55:80:\A);
        % пројекције
        \draw[pomocna] (A) -- (K);
        \draw[pomocna] (A) -- (0,{\A*sin(\t)});
        \draw[pomocna] (K) -- (0,{-\A*sin(\t)});
        \node[left] at (0,{\A*sin(\t)}) {$b$};
        \node[left] at (0,{-\A*sin(\t)}) {$-b$};
        \node[below left] at ({\A*cos(\t)},0) {$a$};
        % углови
        \draw[->] (0.6,0) arc (0:\t:0.6);
        \node at ({\t/2}:0.85) {$\theta$};
        \draw[->, gray!70!black] (0.6,0) arc (0:-\t:0.6);
        \node[gray!70!black] at ({-\t/2}:0.9) {$-\theta$};
        % бројеви
        \draw[->, very thick, blue!70!black] (0,0) -- (A);
        \node[above right, blue!70!black] at (A) {$\underline{A} = a + jb$};
        \node[above left, blue!70!black] at ($(0,0)!0.55!(A)$) {$A$};
        \draw[->, thick, dashed, gray!70!black] (0,0) -- (K);
        \node[below right, gray!70!black] at (K) {$\underline{A}^* = a - jb$};
        \fill[blue!70!black] (A) circle (1.6pt);
        \fill[gray!70!black] (K) circle (1.6pt);
    \end{tikzpicture}
    \caption{Комплексни број $\underline{A}$ и њему конјуговано комплексни број $\underline{A}^*$ у комплексној равни}
    \label{fig:kompleksna-ravan}
\end{figure}

\subsection{Ојлерова формула}

Рекли смо радине да са леве на десно, представљамо шта "види" синус, а са горее на доле шта "види" косинус. Показали смо и да се њихово кретање мође показати да  је само кретање по коодринатним осама, а да је крушно кретање само пројекција та два кретањаа. Показали смо комплексним бројевима да постоје реалан део и имагинаран део који такође образују једну кружницу. 

Спајањем та два добијамо оно што је Ојлер добио пре 300 година. Ојлерова формула повезује експоненцијалну функцију и круг:
\begin{equation}
    e^{j\theta} = \cos\theta + j\sin\theta
\end{equation}
Број $e^{j\theta}$ је тачка на јединичном кругу под углом $\theta$. Када $\theta$ расте, тачка кружи.

За угао $-\theta$ (косинус је паран, а синус непаран):
\[
    e^{-j\theta} = \cos(-\theta) + j\sin(-\theta) = \cos\theta - j\sin\theta
\]
Сабирањем и одузимањем ове две формуле добијамо:
\[
    \cos\theta = \frac{1}{2}\left(e^{j\theta} + e^{-j\theta}\right),
    \qquad
    \sin\theta = \frac{1}{2j}\left(e^{j\theta} - e^{-j\theta}\right)
\]

\begin{kljucno}
\textbf{За запамтити:}
\[
    \underline{A} = A(\cos\theta + j\sin\theta) = A\,e^{j\theta},
    \qquad
    \underline{A}^* = A(\cos\theta - j\sin\theta) = A\,e^{-j\theta}
\]
\end{kljucno}

\begin{figure}[H]
    \centering
    \begin{tikzpicture}[>=Stealth, font=\small]
        \foreach \i/\ang/\natpis in {
            0/0/{e^{j0} = e^{j2\pi} = 1},
            1/90/{e^{j\frac{\pi}{2}} = j},
            2/180/{e^{j\pi} = -1},
            3/270/{e^{j\frac{3\pi}{2}} = e^{-j\frac{\pi}{2}} = -j}} {
            \begin{scope}[xshift=\i*3.6cm]
                \draw[->, gray] (-1.0,0) -- (1.05,0);
                \draw[->, gray] (0,-1.0) -- (0,1.05);
                \draw (0,0) circle (0.75);
                \draw[very thick, blue!70!black] (0,0) -- (\ang:0.75);
                \fill[blue!70!black] (\ang:0.75) circle (2pt);
                \node at (0,-1.4) {$\natpis$};
            \end{scope}
        }
    \end{tikzpicture}
    \caption{Вредности $e^{j\theta}$ за $\theta = 0$, $\frac{\pi}{2}$, $\pi$ и $\frac{3\pi}{2}$}
    \label{fig:euler-primeri}
\end{figure}

Из $e^{j\pi} = -1$ следи \textbf{Ојлеров идентитет}, који повезује бројеве $0$, $1$, $\pi$, $e$ и $j$:
\[
    e^{j\pi} + 1 = 0
\]

\paragraph{Множење са $j$.} Пошто је $j = e^{j\frac{\pi}{2}}$, а при множењу се изложиоци сабирају:
\[
    j\cdot\underline{A} = e^{j\frac{\pi}{2}}\cdot A\,e^{j\theta} = A\,e^{j\left(\theta + \frac{\pi}{2}\right)}
\]
\textbf{Множење са $j$ окреће број за $90^\circ$} супротно од казаљке на сату. Множење са $-j$ окреће га за $90^\circ$ у смеру казаљке.

\paragraph{Веза са наизменичним величинама.} За $\theta = \omega t$ тачка $e^{j\omega t}$ кружи угаоном брзином $\omega$, као навојак генератора. Њен имагинарни део је $\sin\omega t$. Зато наизменичну струју можемо записати као:
\[
    i = I_m\sin(\omega t + \psi) = \operatorname{Im}\left\{ I_m\,e^{j(\omega t + \psi)} \right\}
\]
Вектор који тако кружи у комплексној равни зове се \textbf{фазор}.

\subsection{Примена: отпорник, кондензатор и калем}

Посматрамо струју и напон на три основна елемента. Са $\varphi$ означавамо колико струја \textbf{претходи} напону.

\paragraph{Отпорник.} По Омовом закону $u = R\,i$. Напон и струја истовремено пролазе кроз нулу и максимум. Они су \textbf{у фази}: $\varphi = 0$.

\paragraph{Кондензатор.} Струја је сразмерна брзини промене напона:
\[
    i = C\,\frac{du}{dt}
\]
Ако је $u = U_m\sin\omega t$, онда је (извод синуса је косинус):
\[
    i = \omega C\,U_m\cos\omega t = \omega C\,U_m\sin(\omega t + 90^\circ)
\]
Струја \textbf{претходи} напону за $90^\circ$: $\varphi = +\frac{\pi}{2}$.

\paragraph{Калем.} Напон је сразмеран брзини промене струје:
\[
    u = L\,\frac{di}{dt}
\]
Ако је $i = I_m\sin\omega t$, онда је:
\[
    u = \omega L\,I_m\cos\omega t = \omega L\,I_m\sin(\omega t + 90^\circ)
\]
Напон претходи струји за $90^\circ$. Ако за почетак узмемо напон, $u = U_m\sin\omega t$, онда је $i = I_m\sin(\omega t - 90^\circ)$. Струја \textbf{касни} за напоном за $90^\circ$: $\varphi = -\frac{\pi}{2}$.

\begin{figure}[H]
    \centering
    \begin{tikzpicture}[>=Stealth, font=\small]
        \def\s{0.016667}  % cm по степену (450° = 7,5 cm)
        \def\U{0.85}      % амплитуда напона на цртежу
        \def\I{0.55}      % амплитуда струје на цртежу
        \def\gx{2.6}      % почетак графика
        \def\px{11.9}     % центар фазорског дијаграма
        % --- легенда ---
        \draw[very thick, blue!70!black] ({\gx+2.6},1.6) -- ++(0.5,0) node[right, black] {напон $u$};
        \draw[kriva] ({\gx+4.6},1.6) -- ++(0.5,0) node[right, black] {струја $i$};
        % --- графици и фазори (ред по ред) ---
        \foreach \y/\ph/\ianc in {0/0/above, -2.7/90/left, -5.4/-90/left} {
            \begin{scope}[yshift=\y cm]
                \draw[->] (\gx,0) -- ({\gx+7.9},0) node[right] {$\omega t$};
                \draw[->] (\gx,-1.05) -- (\gx,1.15);
                \draw[very thick, blue!70!black, domain=0:450, samples=100] plot ({\gx+\x*\s},{\U*sin(\x)});
                \draw[kriva, domain=0:450, samples=100] plot ({\gx+\x*\s},{\I*sin(\x+\ph)});
                % фазори
                \draw[->, gray] ({\px-0.4},0) -- ({\px+1.4},0);
                \draw[->, gray] (\px,-1.0) -- (\px,1.0);
                \draw[->, very thick, blue!70!black] (\px,0) -- ({\px+1.15},0) node[below] {$\underline{U}$};
                \draw[->, very thick, red!75!black] (\px,0) -- ++(\ph:0.7) node[\ianc] {$\underline{I}$};
            \end{scope}
        }
        % --- симболи елемената ---
        % отпорник
        \draw (0,0.35) -- (0.45,0.35);
        \draw[thick] (0.45,0.2) rectangle (1.15,0.5);
        \draw (1.15,0.35) -- (1.6,0.35);
        \node at (0.8,-0.15) {Отпорник};
        \node at (0.8,-0.65) {$u = R\,i$};
        % кондензатор
        \begin{scope}[yshift=-2.7cm]
            \draw (0,0.35) -- (0.7,0.35);
            \draw[very thick] (0.7,0.05) -- (0.7,0.65);
            \draw[very thick] (0.9,0.05) -- (0.9,0.65);
            \draw (0.9,0.35) -- (1.6,0.35);
            \node at (0.8,-0.15) {Кондензатор};
            \node at (0.8,-0.75) {$i = C\,\dfrac{du}{dt}$};
            % помак: струја претходи
            \draw[pomocna] ({\gx+360*\s},\I) -- ({\gx+360*\s},1.1);
            \draw[pomocna] ({\gx+450*\s},\U) -- ({\gx+450*\s},1.1);
            \draw[->, thick] ({\gx+450*\s},1.0) -- ({\gx+360*\s},1.0) node[midway, above, font=\footnotesize] {$90^\circ$};
            \draw[->, thick] ({\px+0.35},0) arc (0:90:0.35);
            \node[font=\footnotesize] at ({\px+0.62},0.42) {$90^\circ$};
        \end{scope}
        % калем
        \begin{scope}[yshift=-5.4cm]
            \draw (0,0.35) -- (0.3,0.35);
            \draw[thick] (0.3,0.35) arc (180:0:0.125) arc (180:0:0.125) arc (180:0:0.125) arc (180:0:0.125);
            \draw (1.3,0.35) -- (1.6,0.35);
            \node at (0.8,-0.15) {Калем};
            \node at (0.8,-0.75) {$u = L\,\dfrac{di}{dt}$};
            % помак: струја касни
            \draw[pomocna] ({\gx+90*\s},\U) -- ({\gx+90*\s},1.1);
            \draw[pomocna] ({\gx+180*\s},\I) -- ({\gx+180*\s},1.1);
            \draw[->, thick] ({\gx+90*\s},1.0) -- ({\gx+180*\s},1.0) node[midway, above, font=\footnotesize] {$90^\circ$};
            \draw[->, thick] ({\px+0.35},0) arc (0:-90:0.35);
            \node[font=\footnotesize] at ({\px+0.62},-0.42) {$90^\circ$};
        \end{scope}
    \end{tikzpicture}
    \caption{Напон и струја на отпорнику, кондензатору и калему (лево) и њихови фазори (десно). Код кондензатора струја претходи напону за $90^\circ$, а код калема касни за $90^\circ$.}
    \label{fig:rlc}
\end{figure}

\begin{table}[H]
    \centering
    \renewcommand{\arraystretch}{1.9}
    \begin{tabular}{lccc}
        \toprule
        Елемент & Веза у времену & Комплексни облик & $\varphi$ \\
        \midrule
        отпорник    & $u = R\,i$                 & $\underline{U} = R\,\underline{I}$               & $0$ \\
        кондензатор & $i = C\,\dfrac{du}{dt}$    & $\underline{I} = j\omega C\,\underline{U}$      & $+90^\circ$ \\
        калем       & $u = L\,\dfrac{di}{dt}$    & $\underline{U} = j\omega L\,\underline{I}$      & $-90^\circ$ \\
        \bottomrule
    \end{tabular}
    \caption{Струја и напон на основним елементима}
    \label{tab:rlc}
\end{table}

\paragraph{Шта ту раде комплексни бројеви?} Извод од $e^{j\omega t}$ је $j\omega\,e^{j\omega t}$. Значи, \textbf{извод је множење са $j\omega$}, а множење са $j$ је окрет за $90^\circ$. Зато уместо извода у комплексном облику пишемо само $j\omega$:


\section*{Кратак преглед формула}

\begin{kljucno}
\renewcommand{\arraystretch}{2.3}
\begin{tabular}{ll}
    Тренутна вредност       & $i = I_m\sin(\omega t + \psi)$ \\
    Фреквенција             & $f = \dfrac{1}{T}$ \\
    Кружна фреквенција      & $\omega = \dfrac{2\pi}{T} = 2\pi f$ \\
    Почетна фаза са графика & $\psi = -\dfrac{t_0}{T}\cdot 360^\circ$ \\
    Средња вредност         & $I_{sr} = \dfrac{2}{\pi}\,I_m \approx 0{,}637\,I_m$ \\
    Ефективна вредност      & $I = \dfrac{I_m}{\sqrt{2}} \approx 0{,}707\,I_m$ \\
    Градска мрежа           & $U = 230$~V, \ $f = 50$~Hz, \ $T = 20$~ms, \ $\omega = 314$~rad/s \\
\end{tabular}
\end{kljucno}

\end{document}